\chapter{Выводы}

Выполненная работа была посвящена численному решению модели реального делового цикла для закрытой экономики. В рамках исследования была рассмотрена RBC-модель с репрезентативным домохозяйством и фирмой, выписаны основные условия равновесия, уравнение Эйлера, бюджетное ограничение и соотношения, определяющие стационарное состояние экономики.

Основное внимание в работе уделялось глобальному численному подходу к решению динамической системы. В отличие от локальной линеаризации, которая описывает поведение модели только в малой окрестности стационарного состояния, глобальное решение позволяет анализировать траектории на более широком множестве начальных условий. Это особенно важно для макроэкономических моделей, поскольку при существенных отклонениях от равновесия нелинейные свойства системы начинают играть заметную роль.

В практической части был реализован метод стрельбы, с помощью которого подбиралось начальное потребление, обеспечивающее сходимость траектории к стационарной точке. Для решения неявного уравнения Эйлера на каждом шаге использовался метод Ньютона. Такая схема позволила построить устойчивые переходные траектории капитала, потребления и труда при различных начальных значениях капитала и параметров модели.

Построенные фазовые и трехмерные траектории показали, что экономика сходится к стационарному состоянию как при недостаточном, так и при избыточном начальном уровне капитала. Это подтверждает корректность реализованного численного алгоритма и согласуется с экономической интерпретацией седлово-устойчивой динамики. При этом глобальный подход дает более полное представление о форме переходной траектории, чем локальное приближение.

Сравнение локального и глобального решений показало, что локальные методы удобны для анализа малых отклонений от стационарного состояния, но могут терять точность при больших шоках или далеких начальных условиях. Глобальное решение, напротив, позволяет учитывать особенности модели и ее нелинейности без предварительного сведения системы к линейному виду. Поэтому такой подход является более содержательным при исследовании полной переходной динамики и поведения экономики вне малой окрестности равновесия.

Отдельно следует отметить, что учет нелинейностей открывает возможность анализа эффектов, которые не всегда корректно отражаются в локальной линеаризованной постановке. В частности, в стохастических версиях RBC/DSGE-моделей глобальное решение позволяет выявить эффект сбережений из предосторожности. Его экономический смысл состоит в том, что при наличии неопределенности агенты могут изменять решения о потреблении и накоплении капитала не только из-за среднего ожидаемого шока, но и из-за риска будущих неблагоприятных состояний.

Таким образом, глобальные численные методы дают более гибкий инструмент для анализа макроэкономической динамики. Они позволяют исследовать не только поведение модели около стационарной точки, но и реакцию экономики на существенные отклонения, нелинейные эффекты и потенциальное влияние неопределенности. Полученные результаты могут служить основой для дальнейшего расширения модели, в том числе для учета стохастических шоков, монетарной политики и более детального анализа импульсных откликов.
